\chapter*{Summary}
\label{sommario}
\addcontentsline{toc}{chapter}{Summary}

University bureaucratic information is distributed across official web pages, programme pages, regulations, FAQs, and yearly PDF documents. Students describe their needs in informal language, while institutional sources often use formal administrative terminology. This lexical mismatch can reduce the effectiveness of BM25 retrieval even when the intended answer is present in the corpus.

This thesis investigates whether a low-cost Large Language Model can act as a lightweight query-rewriting layer that improves BM25 retrieval over official UniTrento bureaucratic sources. The LLM is not evaluated as the final answer generator. Instead, it produces a two-layer search query composed of a cleaned version of the student's wording and a small augmentation with institutional terms.

The final benchmark contains 26 Italian student-style questions. Original and rewritten queries are evaluated on two overlapping corpus views: a curated polished corpus and a substantially broader noisy corpus. Three independently generated rewrite sets are frozen and reused unchanged on both corpora, allowing rewrite variability and corpus noise to be analysed separately.

The top three distinct documents are manually labelled as relevant, partially relevant, or irrelevant. Retrieval quality is measured through strict and lenient Hit@1, Hit@3, and Mean Reciprocal Rank. Additional metrics measure whether answerable queries receive useful context beyond the first useful source.

The broader prototype includes controlled source acquisition, manifest and revision tracking, HTML/PDF extraction, OpenSearch BM25 retrieval, backend APIs, a frontend interface, and debugging tools. These components provide the product context, while the scientific core remains a controlled and reproducible evaluation of query rewriting as a thin semantic layer over deterministic retrieval.
\par
